%% LyX 2.4.4 created this file.  For more info, see https://www.lyx.org/.
%% Do not edit unless you really know what you are doing.
\documentclass[english]{article}
\usepackage[T1]{fontenc}
\usepackage[latin9]{inputenc}
\usepackage{enumitem}
\usepackage{amsmath}
\usepackage{amsthm}
\usepackage{geometry}
\geometry{verbose,tmargin=1in,bmargin=1in,lmargin=1in,rmargin=1in}
\PassOptionsToPackage{normalem}{ulem}
\usepackage{ulem}

\makeatletter
%%%%%%%%%%%%%%%%%%%%%%%%%%%%%% Textclass specific LaTeX commands.
\numberwithin{equation}{section}
\numberwithin{figure}{section}
\newlength{\lyxlabelwidth}      % auxiliary length 

%%%%%%%%%%%%%%%%%%%%%%%%%%%%%% User specified LaTeX commands.
\usepackage{fancyhdr}
\usepackage{lastpage}
\pagestyle{fancy}
\usepackage{enumitem}
\usepackage{ifthen}
\everymath=\expandafter{\the\everymath\displaystyle}
%\DeclareMathSizes{10}{10}{10}{10}
\usepackage{multicol}

\usepackage{tikz}
\usepackage{tikz-3dplot}
\usetikzlibrary{calc,arrows}

\usetikzlibrary{patterns}
\usetikzlibrary{plotmarks}
\usepackage{pgfplots}
\pgfplotsset{compat=newest}

\usepackage{calc}
\usepackage[nomessages]{fp}

\usepackage{datetime}
\usepackage[super]{nth}
% Added by lyx2lyx
\setlength{\parskip}{\smallskipamount}
\setlength{\parindent}{0pt}

\makeatother

\usepackage{babel}
\begin{document}
\renewcommand{\author}{Dr.\,James Ashe}

\newcommand{\classNum}{336}

\newcommand{\classSec}{A}

\newcommand{\examType}{Exam 3 Review}

\renewcommand{\date}{\formatdate{18}{11}{2013}}

%%First page's header%%%%%%%%%%%%%%%%%%%%%%
\fancypagestyle{firststyle} %%header settings for first pagr
{
	\fancyhead[L]{MTH \classNum{}(\classSec{}) \author{}\\
	\textbf{\examType{}} ~\date{}}
	\fancyhead[C]{}
	\fancyhead[R]{\textbf{}}
	\setlength{\headsep}{14pt}
}
%%%%%%%%%%%%%%%%%%%%%%%%%%%%%%%%%%%%%%%%%%

%%Next page's header%%%%%%%%%%%%%%%%%%%%%%%
\setlist{nolistsep}
\lhead{\classNum{}(\classSec{})} 
\chead{\textbf{}} 
\cfoot{\thepage{}~of~\pageref{LastPage}} 
\setlength{\headsep}{5pt} 
%%%%%%%%%%%%%%%%%%%%%%%%%%%%%%%%%%%%%%%%%%%

%%Various settings; see the preamble for other settings%%%%%
%\setenumerate[0]{leftmargin=12pt,itemindent=0pt,labelwidth=0pt}
\setlist{topsep=.5pt}
\renewcommand{\labelenumi}{\textbf{\arabic{enumi}.}}
\renewcommand{\labelenumii}{\textbf{\alph{enumii}.}}
\thispagestyle{firststyle}
%%%%%%%%%%%%%%%%%%%%%%%%%%%%%%%%%%%%%%%%%%

\newcommand{\xmax}{0}
\newcommand{\xmin}{-\xmax}
\newcommand{\xstep}{0}
\newcommand{\ymax}{0}
\newcommand{\ymin}{-\ymax}
\newcommand{\ystep}{0} 
\newcommand{\dspace}{\vspace{1cm}}
\newcommand{\xa}{0}
\newcommand{\xb}{0}
\newcommand{\ya}{1}
\newcommand{\yb}{1}

\subsubsection*{Chapter 5 Determinants}
\begin{itemize}
\item \textbf{5.1 Properties of Determinants:}

\begin{itemize}
\item Calculate the determinant of a matrix using properties of determinants.

\begin{itemize}
\item [\textbf{ex.}] HW 1ab.
\item [\textbf{ex.}]Let $\det(A)=\sqrt{2}$. Does $A^{-1}$ exist? If so
find $\det(A^{-1})$.
\item [\textbf{ex.}]Let $R$ be the reduced row echelon form of $A$. How
does $\det A$ relate to $\det R$?
\item [\textbf{ex.}]Let $A\in\mathcal{M}_{2\times2}$ with a determinant
of $3$ and $B=\begin{bmatrix}\begin{array}{rr}
1 & 2\\
2 & 2
\end{array}\end{bmatrix}$, find $\det AB$.
\end{itemize}
\item For $n\times n$ matrices:

\begin{itemize}
\item See Proposition 1.7.
\item $A$ is an upper (or lower) triangle matrix $\Rightarrow$ $\det A$
is equal to the product of $A$'s diagonal entries.
\item $A$ is nonsingular $\iff$ $\det A\neq0$.
\item $\det(AB)=\det A\det B$.
\item $A$ is a nonsingular $\Rightarrow$$\det(A^{-1})=\frac{1}{\det A}$.
\end{itemize}
\end{itemize}
\item \textbf{5.2 Cofactors}

\begin{itemize}
\item Calculate the determinant of a matrix using cofactors.
\item Calculate the $x_{i}$ coordinate in a solution vector using Cramer's
rule.
\end{itemize}
\end{itemize}

\subsubsection*{Chapter 6 Eigenvalues and Eigenvectors}
\begin{itemize}
\item \textbf{6.1: The Characteristic Polynomial}

\begin{itemize}
\item Calculate the characterisitice polynomial of $A$, $p_{A}(t)$.
\item Find the eigenvalues and eigenvectors of $A$.
\item Know the definition of \emph{eigenvalue} and \emph{eigenvector}.
\end{itemize}
\item \textbf{6.2: Diagonalization}

\begin{itemize}
\item Know the defintion of \emph{similar matrices}: $A$ and $B$ are similar
if there exists an invertible matrix $P$ such that $A=P^{-1}BP$
\item \uline{Prove:} If $A$ and $B$ are similar matrices then $\det(A)=\det(B)$.
(Hint: use the last two properties under 5.1 above)
\item Know the definition of \emph{diagonalizable}.
\end{itemize}
\end{itemize}

\end{document}
